\input{config/documento}

%Metadados
\titulo{Título do trabalho}
\autor{Nome 1 \and \\
Nome 2 \and \\
Nome 3}
\newcommand\disciplina{Estatística}
\newcommand\professor{Aléssio Tony C. Almeida}


% ----------------------------------------------------------
% Início do documento
% ----------------------------------------------------------
\begin{document}

\input{config/pretextual}

\begin{resumo}
Este trabalho tem por objetivo \lipsum[1-1]

\textbf{Palavras-chave}: Palavra 1. Palavra 2. Palavra 3. Palavra 4. Palavra 5.
\end{resumo}

\textual
\parindentABNT

\chapter{Introdução}

\lipsum[1-1]

%Exemplos de citação:

\begin{itemize}
    \item De acordo com \citeonline{Almeida2019} e \citeonline{Silva2019}, o Brasil...
    \item O Programa Farmácia Popular do Brasil visa aumentar o acesso a medicamentos para doentes crônicos no país, usando especialmente a rede conveniada ao programa \cite{Almeida2019}.
\end{itemize}

%Texto aleatório
\lipsum[1-5]

\chapter{Metodologia}

Descrever nesta parte a metodologia do seu trabalho.

\section{Modelo empírico}

Apresentar o modelo empírico utilizado.

\section{Dados}

Descrever as variáveis utilizadas, período de análise, abrangência territorial do estudo, tamanho da amostra e as fontes de dados, bem como apresentar uma tabela com uma sumarização estatística das variáveis usadas (média, desvio-padrão, mínimo, máximo).

\chapter{Resultados}

Apresentar os achados da pesquisa, bem como discussões acerca das principais implicações desses resultados encontrados.

\chapter{Conclusões}

Informar neste capítulo as principais considerações da pesquisa: resultados-chave, diferenciais do trabalho, limitações, perspectivas futuras de trabalhos relacionados ao tema.

\postextual
\bibliography{referencias.bib}

\end{document}
